% ═══════════════════════════════════════════════════════════════
% MUSTERLÖSUNG: HAY (Repaso) + ESTAR
% ═══════════════════════════════════════════════════════════════
% Thema: Hay vs Estar
% Klasse 7 | Unidad 2 | Mi instituto | DS 02
% ═══════════════════════════════════════════════════════════════

\documentclass[10pt, a4paper]{article}

\usepackage[utf8]{inputenc}
\usepackage[T1]{fontenc}
\usepackage[ngerman]{babel}

% --- SCHRIFTEN ---
\usepackage{palatino}
\usepackage{helvet}
\newcommand{\sanstext}[1]{{\fontfamily{phv}\selectfont #1}}

\usepackage{microtype}
\usepackage[margin=1.4cm, top=1.6cm, bottom=1.3cm]{geometry}
\usepackage{enumitem}
\usepackage{tabularx}
\usepackage{booktabs}
\usepackage{array}
\usepackage{multicol}
\usepackage{tikz}
\usetikzlibrary{calc}

\usepackage{fancyhdr}
\usepackage{lastpage}
\usepackage{needspace}

\usepackage{xcolor}
\usepackage{amssymb}
\usepackage{tcolorbox}
\tcbuselibrary{skins}

% ═══════════════════════════════════════════════════════════════
% FARBEN
% ═══════════════════════════════════════════════════════════════

\definecolor{kopfzeile}{HTML}{1A1A1A}
\definecolor{akzent}{HTML}{C41E3A}          % Spanisch-Karminrot
\definecolor{hellgrau}{HTML}{F2F2F2}
\definecolor{mittelgrau}{HTML}{777777}
\definecolor{dunkelgrau}{HTML}{444444}
\definecolor{loesung}{HTML}{C41E3A}         % Rot für Lösungen

% ═══════════════════════════════════════════════════════════════
% VARIABLEN
% ═══════════════════════════════════════════════════════════════

\newcommand{\thema}{Hay y estar}
\newcommand{\klasse}{7}
\newcommand{\maxpunkte}{34}

% ═══════════════════════════════════════════════════════════════
% KOPF- UND FUSSZEILE
% ═══════════════════════════════════════════════════════════════

\pagestyle{fancy}
\fancyhf{}
\renewcommand{\headrulewidth}{0pt}
\renewcommand{\footrulewidth}{0pt}
\cfoot{\sanstext{\footnotesize\textcolor{mittelgrau}{\thepage\,/\,\pageref{LastPage}}}}

% ═══════════════════════════════════════════════════════════════
% BOXEN
% ═══════════════════════════════════════════════════════════════

\newtcolorbox{grammatikbox}{
    colback=hellgrau,
    colframe=hellgrau,
    boxrule=0pt,
    arc=0pt,
    left=8pt, right=8pt, top=8pt, bottom=6pt,
    borderline north={1.5pt}{0pt}{dunkelgrau}
}

\newtcolorbox{merkbox}{
    colback=hellgrau,
    colframe=hellgrau,
    boxrule=0pt,
    arc=0pt,
    left=8pt, right=8pt, top=8pt, bottom=6pt,
    borderline north={1.5pt}{0pt}{akzent}
}

\newcommand{\regel}[1]{%
    \tikz[baseline=(X.base)]\node[fill=akzent, text=white,
    font=\sanstext\scriptsize\bfseries, inner sep=2pt, rounded corners=1pt](X){#1};%
}

% ═══════════════════════════════════════════════════════════════
% BEFEHLE
% ═══════════════════════════════════════════════════════════════

\newcommand{\es}[1]{\textbf{#1}}
\newcommand{\de}[1]{{\small\textit{\textcolor{mittelgrau}{#1}}}}
\newcommand{\luecke}[1]{\rule{#1}{0.4pt}}
\newcommand{\punktebox}[1]{\hfill\sanstext{\small[\_\_/#1]}}

% LÖSUNG in Rot
\newcommand{\lsg}[1]{\textcolor{loesung}{\textbf{#1}}}

\newcommand{\aufgabe}[2]{%
    \needspace{4\baselineskip}%
    \vspace{10pt}%
    \noindent\sanstext{\textbf{#1}\quad#2}%
    \vspace{5pt}%
    \nopagebreak[4]%
}

\newlist{teil}{enumerate}{1}
\setlist[teil]{label=\alph*), leftmargin=1.8em, itemsep=4pt, parsep=0pt, topsep=3pt}

\setlength{\parindent}{0pt}
\setlength{\parskip}{4pt}
\setlength{\columnsep}{24pt}

% ═══════════════════════════════════════════════════════════════
% DOKUMENT
% ═══════════════════════════════════════════════════════════════

\begin{document}

% ═══════════════════════════════════════════════════════════════
% HEADER
% ═══════════════════════════════════════════════════════════════

\noindent
\begin{tikzpicture}[remember picture, overlay]
    \fill[kopfzeile] (current page.north west) rectangle +(\paperwidth, -1cm);
    \node[anchor=west, text=white, font=\sanstext\large\bfseries]
        at ([xshift=1.4cm, yshift=-0.5cm]current page.north west) {\thema{} -- \textcolor{yellow}{MUSTERLÖSUNG}};
    \node[anchor=east, text=white, font=\sanstext\small]
        at ([xshift=-1.4cm, yshift=-0.5cm]current page.north east)
        {Spanisch\,·\,Klasse \klasse\,·\,Arbeitsblatt};
\end{tikzpicture}

\vspace{0.5cm}

% --- Name/Datum ---
\noindent
\sanstext{\small\textbf{Name:}} \lsg{MUSTERLÖSUNG} \hfill
\sanstext{\small\textbf{Datum:}} \luecke{2.5cm}

\vspace{8pt}

% ═══════════════════════════════════════════════════════════════
% GRAMMATIK-ÜBERSICHT
% ═══════════════════════════════════════════════════════════════

\begin{grammatikbox}
\begin{multicols}{2}

{\large\bfseries hay} \de{es gibt}

\vspace{3pt}

{\small
\begin{tabular}{@{}ll@{}}
+ un/una & \es{Hay una mesa.} \\
+ Zahl & \es{Hay tres libros.} \\
+ mucho/poco & \es{Hay muchos cuadernos.} \\
+ ohne Artikel & \es{Hay libros.} \\
\end{tabular}}

\regel{KEIN el/la nach hay!}

\columnbreak

{\large\bfseries estar} \de{sein (Ort/Zustand)}

\vspace{3pt}

{\small
\begin{tabular}{@{}ll@{}}
yo & \es{estoy} \\
tú & \es{estás} \\
él/ella & \es{está} \\
\end{tabular}
\hspace{0.5cm}
\begin{tabular}{@{}ll@{}}
nosotros & \es{estamos} \\
vosotros & \es{estáis} \\
ellos & \es{están} \\
\end{tabular}}

\regel{+ el/la (bestimmt)}

\end{multicols}
\end{grammatikbox}

\vspace{2pt}

% ═══════════════════════════════════════════════════════════════
% AUFGABEN MIT LÖSUNGEN
% ═══════════════════════════════════════════════════════════════

\begin{multicols}{2}

% --- AUFGABE 1 ---
\aufgabe{1}{Wann benutzt man \es{hay}? Ordne zu.} \punktebox{4}

{\small Verbinde mit Linien.}

\vspace{6pt}

\begin{tikzpicture}[
    regelbox/.style={draw, rectangle, minimum width=2.6cm, minimum height=1.1cm, font=\small, align=center},
    beispielbox/.style={draw, rectangle, minimum width=3.2cm, minimum height=0.65cm, font=\small}
]
% Linke Spalte: Regeln (mit deutscher Erklärung)
\node[regelbox] (R1) at (0,0) {hay + un/una\\[-1pt]{\tiny\textcolor{mittelgrau}{unbest. Artikel}}};
\node[regelbox] (R2) at (0,-1.3) {hay + número\\[-1pt]{\tiny\textcolor{mittelgrau}{Zahl}}};
\node[regelbox] (R3) at (0,-2.6) {hay + mucho/poco\\[-1pt]{\tiny\textcolor{mittelgrau}{Mengenangabe}}};
\node[regelbox] (R4) at (0,-3.9) {hay + sin artículo\\[-1pt]{\tiny\textcolor{mittelgrau}{ohne Artikel}}};

% Rechte Spalte: Beispiele (gemischt)
\node[beispielbox] (B1) at (4.8,0) {Hay muchos libros.};
\node[beispielbox] (B2) at (4.8,-1.3) {Hay libros.};
\node[beispielbox] (B3) at (4.8,-2.6) {Hay una mesa.};
\node[beispielbox] (B4) at (4.8,-3.9) {Hay tres sillas.};

% Lösungslinien in Rot
\draw[loesung, thick] (R1.east) -- (B3.west);  % un/una → Hay una mesa
\draw[loesung, thick] (R2.east) -- (B4.west);  % número → Hay tres sillas
\draw[loesung, thick] (R3.east) -- (B1.west);  % mucho/poco → Hay muchos libros
\draw[loesung, thick] (R4.east) -- (B2.west);  % sin artículo → Hay libros
\end{tikzpicture}

\vspace{6pt}

% --- AUFGABE 2 ---
\aufgabe{2}{Bilde Sätze mit \es{hay}.} \punktebox{6}

{\small Benutze verschiedene Strukturen (un/una, Zahl, mucho, ohne Artikel).}

\vspace{4pt}

\begin{teil}
    \item (una / pizarra) En la clase \lsg{hay una pizarra.}

    \vspace{2pt}

    \item (25 / sillas) En la clase \lsg{hay 25 sillas.}

    \vspace{2pt}

    \item (muchos / bolígrafos) En mi estuche \lsg{hay muchos bolígrafos.}

    \vspace{2pt}

    \item (libros) En la mochila \lsg{hay libros.}

    \vspace{2pt}

    \item (pocos / cuadernos) En mi mesa \lsg{hay pocos cuadernos.}

    \vspace{2pt}

    \item (un / ordenador) En la clase \lsg{hay un ordenador.}
\end{teil}

\columnbreak

% --- AUFGABE 3 ---
\aufgabe{3}{Ergänze \es{estar} und die Präposition.} \punktebox{6}

{\small Benutze: \es{en, encima de, debajo de, al lado de, delante de, detrás de}}

\vspace{4pt}

{\footnotesize\textit{Beispiel:} El libro / la mesa $\rightarrow$ \es{El libro está encima de la mesa.}}

\vspace{4pt}

\begin{teil}
    \item ¿Dónde \lsg{está} el bolígrafo?

    \vspace{2pt}

    \item El bolígrafo \lsg{está} \lsg{encima de} la mesa.

    \vspace{2pt}

    \item La mochila \lsg{está} \lsg{debajo de} la silla.

    \vspace{2pt}

    \item El profesor \lsg{está} \lsg{delante de} la pizarra.

    \vspace{2pt}

    \item Los libros \lsg{están} \lsg{en} la mochila.

    \vspace{2pt}

    \item Yo \lsg{estoy} \lsg{al lado de} la ventana.
\end{teil}

\vspace{6pt}

% --- AUFGABE 4 ---
\aufgabe{4}{¿\es{Hay} o \es{estar}? Elige.} \punktebox{6}

{\small Kreuze an und ergänze die richtige Form.}

\vspace{4pt}

\begin{teil}
    \item \lsg{Hay} un libro en mi mochila.

    {\footnotesize \lsg{$\boxtimes$ hay} \hspace{1cm} $\square$ estar}

    \vspace{3pt}

    \item ¿Dónde \lsg{está} la goma?

    {\footnotesize $\square$ hay \hspace{1cm} \lsg{$\boxtimes$ estar: está}}

    \vspace{3pt}

    \item En la clase \lsg{hay} muchas ventanas.

    {\footnotesize \lsg{$\boxtimes$ hay} \hspace{1cm} $\square$ estar}

    \vspace{3pt}

    \item El cuaderno \lsg{está} encima de la mesa.

    {\footnotesize $\square$ hay \hspace{1cm} \lsg{$\boxtimes$ estar: está}}

    \vspace{3pt}

    \item ¿\lsg{Hay} lápices en tu estuche?

    {\footnotesize \lsg{$\boxtimes$ hay} \hspace{1cm} $\square$ estar}

    \vspace{3pt}

    \item Mi mochila \lsg{está} debajo de la silla.

    {\footnotesize $\square$ hay \hspace{1cm} \lsg{$\boxtimes$ estar: está}}
\end{teil}

\end{multicols}

% ═══════════════════════════════════════════════════════════════
% AUFGABE 5: Volle Breite
% ═══════════════════════════════════════════════════════════════

\vspace{2pt}
\noindent\textcolor{hellgrau}{\rule{\textwidth}{1.5pt}}
\vspace{2pt}

\aufgabe{5}{Beschreibe deinen Klassenraum.} \punktebox{8}

{\small Schreibe 4 Sätze mit \es{hay} und 4 Sätze mit \es{estar}. Benutze verschiedene Präpositionen!}

\vspace{4pt}

\begin{multicols}{2}

\textbf{Con \es{hay}:}

\vspace{3pt}

1. \lsg{En mi clase hay una pizarra grande.}

\vspace{4pt}

2. \lsg{Hay veinticinco mesas.}

\vspace{4pt}

3. \lsg{Hay muchos libros en la estantería.}

\vspace{4pt}

4. \lsg{Hay un ordenador.}

\columnbreak

\textbf{Con \es{estar}:}

\vspace{3pt}

1. \lsg{El ordenador está encima de la mesa.}

\vspace{4pt}

2. \lsg{La papelera está al lado de la puerta.}

\vspace{4pt}

3. \lsg{Mi mochila está debajo de mi silla.}

\vspace{4pt}

4. \lsg{Los mapas están en la pared.}

\end{multicols}

\vspace{4pt}

% ═══════════════════════════════════════════════════════════════
% AUFGABE 6: Bonus
% ═══════════════════════════════════════════════════════════════

\aufgabe{Bonus}{Mini-Diálogo: ¿Dónde está...?} \punktebox{4}

{\small Schreibe einen Dialog (4-6 Zeilen). Person A sucht etwas, Person B hilft.}

\vspace{4pt}

{\footnotesize\textit{Beispiel:} A: ¿Hay un bolígrafo? -- B: Sí, hay uno. -- A: ¿Dónde está? -- B: Está encima de la mesa.}

\vspace{4pt}

\lsg{Beispiel-Dialog:}

\lsg{A: ¿Hay una goma en tu estuche?}

\lsg{B: Sí, hay una.}

\lsg{A: ¿Dónde está?}

\lsg{B: Está en mi estuche, encima de la mesa.}

\lsg{A: Gracias.}

\lsg{B: De nada.}

% ═══════════════════════════════════════════════════════════════
% MERKKASTEN
% ═══════════════════════════════════════════════════════════════

\vspace{6pt}

\begin{merkbox}
\sanstext{\textbf{Merke: hay vs estar}}

\vspace{3pt}

\begin{multicols}{2}

{\small
\textbf{HAY} = \textit{Es gibt} (Existenz)
\begin{itemize}[leftmargin=1em, itemsep=0pt]
    \item + un/una, Zahl, mucho/poco, ohne Artikel
    \item Frage: \es{¿Qué hay...? ¿Cuántos hay?}
\end{itemize}}

\columnbreak

{\small
\textbf{ESTAR} = \textit{Wo ist? Wie ist?} (Ort/Zustand)
\begin{itemize}[leftmargin=1em, itemsep=0pt]
    \item + el/la, mi/tu (bestimmt)
    \item Frage: \es{¿Dónde está...?}
\end{itemize}}

\end{multicols}
\end{merkbox}

% ═══════════════════════════════════════════════════════════════
% FOOTER
% ═══════════════════════════════════════════════════════════════

\vfill

\noindent\rule{\textwidth}{0.5pt}

\vspace{4pt}

\hfill\sanstext{\textbf{Punkte:}} \lsg{34} / \maxpunkte

\end{document}
